% Begin text transcription and conversion to LaTeX

\chapter{Principios básicos de protección contra radiación externa}

Son tres los factores más importantes que deben tomarse en cuenta a fin de que la dosis que se reciba sea mínima: Tiempo, distancia y blindaje.

Mantener el menor tiempo a la mayor distancia posible he interponiendo el mayor blindaje posible entre la fuente de radiación y la persona, esto garantiza que la dosis será mínima.

La exposición sigue la ley del cuadrado inverso, esto es, varía en razón inversa al cuadrado de la distancia. Matemáticamente se expresa como:
\[
X_1 r_1^2 = X_2 r_2^2
\]
donde \(X_1\) es la rapidez de exposición a la distancia \(r_1\) y \(X_2\) es la rapidez de exposición a la distancia \(r_2\).

Así, si se desea calcular la rapidez de exposición a 4 metros de una cierta fuente puntual que esta dando \(10 \text{ R/h}\) a un metro de distancia se procedería como sigue:
\[
X_2 = \frac{(r_1/r_2)^2 X_1}{1} = \frac{(1/4)^2 \cdot 10}{1} = 0.0625 \cdot 10 = 0.625 \text{ R/hr}
\]
\[
X_2 = 625 \text{ mR/hr}
\]

Como se observa, alejarse de 1 a 4 metros reduce la rapidez de exposición de \(10 \text{ R/hr}\) a \(0.625 \text{ mR/hr}\).

Por otra parte, el blindaje es un método muy eficiente para reducir la rapidez de exposición. Los rayos gamma se atenúan al pasar por un material y el grado de atenuación depende del material en cuestión: por ejemplo se requieren \(6.2\) pulgadas de concreto para reducir la intensidad de radiación a un décimo de su valor original cuando la fuente es Ir-192, mientras que solo se requieren \(0.64\) pulgadas de plomo para reducirla en la misma proporción.

De esto se deduce que el plomo es mejor blindaje que el concreto.

Al espesor necesario de un cierto material que reduce la intensidad de radiación a la mitad de su valor original se le llama "capa hemirreductora" y es distinta para cada radioisótopo y para cada material.

% End of LaTeX content transcription
